% ============================ 5. Results ================================
\section{Results}\label{sec:results}

All numbers in this section are frozen archive results
(Section~\ref{sec:methods}); each claim is traceable to a recorded
experiment.

\subsection{History dependence and non-Markovianity}\label{sec:res-phase1}

Natural collisions were abundant for every window-carrying model
(found pairs: B1 92{,}600; B2 19{,}234; B3 1{,}633{,}844; B4
1{,}427{,}817; 1{,}500 analysed per model; Figure~\ref{fig:phase1},
Table~\ref{tab:phase1}).  On the continuous branch, the mean one-step
amplification $K$ was $0.5000000$ for B1 (exact; maximum deviation
$1.8\times10^{-8}$, machine precision), $4.47\times10^{3}$ (95\% CI
$[3.26,5.92]\times10^{3}$) for B2, $1.235\times10^{4}$
($[7.78,1.83]\times10^{4}$) for B3 and $2.047\times10^{4}$
($[1.12,3.57]\times10^{4}$) for B4, with $73.7\%$ of B4 collision pairs
exceeding $K > 10^{3}$.  Post-reset Markov violations
($|\Delta h_{10}| > \delta$) occurred in 0\%, 0.6\%, 33.1\% and 32.2\%
of pairs for B1--B4 respectively.  When the probe was not surprising
relative to the local window (gates closed, $S_{10} = 0$), the current
vanished and $K = 0.5$ exactly for both B3 and B4: surprise-gating makes
the non-Markovianity \emph{conditional} on prediction error.  Attention
divergence predicted the current divergence (Pearson
$r(JSD, \Delta A) = 0.87$--$0.89$, $p < 10^{-300}$; $\Delta A$ vs.\
$\Delta\tilde{h}$ $r = 1.000$), and the amplification ordering was
B4 $>$ B3 $>$ B2 (median $K$ ratios 2.15 and 3.29, Mann--Whitney
$p = 1.3\times10^{-32}$ and $p = 1.3\times10^{-63}$).  After a common
double spike at $T$ (843 fully resynchronised pairs), B1 stayed
synchronised forever (machine-exact), B2 re-diverged on the continuous
branch and re-converged after the window flush, while B3/B4 retained
inter-pair gaps beyond the flush in $\sim$25\% of pairs: a spike reset
erases $h$ but not the windowed memory.  These results establish that
the scalar membrane potential is not a sufficient state whenever the
update reads the windowed past, for \emph{all} window-carrying models;
history dependence is not an attention-specific effect.

\subsection{Local fluctuation geometry}\label{sec:res-phase2}

With the shared four-dimensional frame $z_C$ and a matched full frame
$z_F$ (Section~\ref{sec:methods-phase2}; Figure~\ref{fig:geometry},
Figure~\ref{fig:rank}, Table~\ref{tab:phase2}), event-locked velocity
effective rank $d_D$ (mean over $\tau \in \{0,1,2\}$, 374 pooled
events) was: B1 $1.674$ ($\Delta$ vs.\ quiet $0.674$), B3 $1.458$
($z_C$) and $1.484$ ($z_F$), B4 $1.589$ ($z_C$) and $2.089$ ($z_F$;
$\Delta = 1.089$), and the input-only control C0 $1.055$.  Three
findings follow.  First, event-locked expansion exists for every
h-bearing model and is not specific to TAN: in the shared frame B4 does
not exceed B1 ($1.589 < 1.674$), and this negative contrast is robust
across amplitude controls (raw/stratified/residualised
$1.589/1.701/1.712$ vs.\ B1 $1.674/1.732/1.747$).  Second, within the
matched full frame that contains the attention coordinates, B4 exceeds
B3 in effective rank (event $d_D$ 2.089 vs.\ 1.484; event CIs
$[1.87,2.30]$ vs.\ $[1.42,1.53]$, disjoint) and in covariance scale
(event $\log\operatorname{Tr}$ 3.58 vs.\ 2.48 nats; $\Delta = +1.09$
nats); the rank contrast survives and slightly strengthens under
amplitude controls ($+0.61$ raw, $+0.74$ stratified, $+0.78$
residualised).  Third, the velocity-geometry response is temporally
structured: its peak occurs at $\tau \approx 4$--$6$ (content-rich
windows after burst onsets, $d_D$ up to 2.89), not at the first pulse
($d_D(\tau{=}0) \approx 1.66$), and recovery is faster for B4
($\tau_{e\text{-fold}} = 5.4$, CI $[4.6,6.3]$, vs.\ B3 8.5,
$[6.1,10.6]$).  These are statements about the effective rank of local
response covariance; they are not statements about intrinsic dimension.

\subsection{Target-relevant temporal retrieval (Probe 1)}\label{sec:res-probe1}

Pooled balanced three-class accuracy over the frozen three seeds
(Figure~\ref{fig:probe1}, Table~\ref{tab:probe1}) was, for clean /
distractor conditions: B1 $0.331 / 0.340$ (chance), B2
$0.867 / 0.843$ (strong, stable positional-memory baseline), B3
$0.375 / 0.337$ (with lo-class-blind cells; one cell invalid for
interpretation), B4 $0.328 / 0.389$.  The B4 vs.\ B3 contrast in the
distractor condition was $+0.0570$ (95\% CI $[+0.0387, +0.0758]$,
$n = 6000$ pooled trials), with per-seed differences $+0.056$, $+0.043$,
$+0.056$.  B4's absolute elevation over chance in the distractor
condition attenuated across seeds ($0.427$, $0.398$, $0.342$;
seed 20260906 is $\approx$ chance).  Taken together, and with B3's
readout limitation and the task's structural bias toward fixed-position
memory (B2 dominates), the evidence supports the constrained statement:
Full TAN shows evidence of improved target-relevant retrieval under
content-rich temporal context relative to the no-attention control, with
a partially replicated, seed-attenuated effect.  This is \emph{limited
support}, not a demonstration of flexible temporal retrieval or
distractor robustness.

\subsection{Identifiability audit of query-conditioned binding (Probe 2)}
\label{sec:res-probe2}

The identifiability audit did not support genuine query-conditioned
routing in the current scalar TAN kernel (Figure~\ref{fig:probe2},
Table~\ref{tab:probe2}).  The B2 linear shortcut test reached balanced
accuracy $0.779$ ($[0.724, 0.836]$) and the B3 shortcut test $0.709$
($[0.644, 0.775]$) against a red line of 0.58, so the task as encoded
was solvable through storage/integration shortcuts before any routing
mechanism was evaluated.  For B4, the attention hit rate (argmax of the
history-tap attention weights landing on the target position) was $0.511$
(bootstrap CI lower bound $0.481$), split as $0.000$ for $q = A$ and
$1.000$ for $q = B$; the target/distractor contrast was $R = 0.034$
(required $\geq 2$); and the query-conditioned JSD of attention
distributions was $0.0002$ with permutation $p = 0.92$ (no detectable
query dependence).  The decisive counterfactual: with history held
fixed and only the query changed ($q = A \to B$), the attention argmax
flip rate was $0.000$.  Changing the query changed attention sharpness,
never the winner among history taps.

\subsection{Mechanistic diagnosis}\label{sec:res-mechanism}

Three layers explain the audit outcome (Figure~\ref{fig:mechanism}).
\textit{Layer 1, encoding shortcut:} the composite encoding used
asymmetric key amplitudes, so B-coded events were systematically larger;
storage/integration readouts could exploit static amplitude, and B4's
amplitude-based attention was attracted to B events ($q=B$ hit 1.000,
$q=A$ hit 0.000).  \textit{Layer 2, scalar-kernel limitation:}
independently of the encoding, with $E_{t,i} = \gamma_t x_i$ the ratio
$E_{t,i}/E_{t,j} = x_i/x_j$ is query-independent for fixed history;
query identity enters only through the scalar $\gamma_t$, so query
changes weighting strength, not selected content.  \textit{Layer 3,
self-attraction:} query amplitudes (3.0/4.0) exceeded all history codes,
so the softmax concentrated on the current tap in 100\% of trials -- an
experimental confound that does not explain away the structural
limitation of the scalar kernel.
